\label{ejercicios}
\chapter{Ejercicios y estudios}

\newcounter{exercise}

Este capítulo presenta una serie de ejercicios con tres objetivos: desarrollar la técnica, interiorizar algunos de los patrones rítmicos más importantes de la música árabe, y afianzar el conocimiento de los distintos \maqamaatb.

Se recomienda practicarlos a velocidad moderada, haciendo especial énfasis en lograr una interpretación clara y una buena pulsación, para después ir trabajando a un \emph{tempo} más elevado.

\section{Estudios y ejercicios progresivos}

Los siguientes estudios y ejercicios se han adaptado de \emph{Método de oud} de Yamil Bashir, y corresponden a los dos primeros años de estudio del ud.

\newexercise{}


\begin{center}
\begin{lilypond}
\include "arabic.ly"
notes = \relative {
do,1 fa la re sol do
\bar "||"
}

\layout {
  \omit Voice.StringNumber
  \context {
      \TabStaff
      \tabFullNotation
      stringTunings = \stringTuning <do, fa, la, re sol do'>
    }
  indent = #0
  \context {
    \Score
    supportNonIntegerFret = ##t
  }
}

<<
  \new TabStaff <<\notes>>
>>
\end{lilypond}
\end{center}


\newexercise{}


\begin{center}
\begin{lilypond}
\include "arabic.ly"
notes = \relative {
fa,1 la re sol do do sol re la fa
\bar "||"
}

\layout {
  \omit Voice.StringNumber
  \context {
      \TabStaff
      \tabFullNotation
      stringTunings = \stringTuning <do, fa, la, re sol do'>
    }
  indent = #0
  \context {
    \Score
    supportNonIntegerFret = ##t
  }
}

<<
  \new TabStaff <<\notes>>
>>
\end{lilypond}
\end{center}


\newexercise{}


\begin{center}
\begin{lilypond}
\include "arabic.ly"
notes = \relative {
fa,2 la re sol do sol re la fa sol' fa,1
\bar "||"
}

\layout {
  \omit Voice.StringNumber
  \context {
      \TabStaff
      \tabFullNotation
      stringTunings = \stringTuning <do, fa, la, re sol do'>
    }
  indent = #0
  \context {
    \Score
    supportNonIntegerFret = ##t
  }
}

<<
  \new TabStaff <<\notes>>
>>
\end{lilypond}
\end{center}


\newexercise{}


\begin{center}
\begin{lilypond}
\include "arabic.ly"
notes = \relative {
fa,2 la fa re' re sol sol re sol do sol re la re fa,1 fa2 la fa la re sol re sol sol do sol do do sol re la fa sol' sol1
\bar "||"
}

\layout {
  \omit Voice.StringNumber
  \context {
      \TabStaff
      \tabFullNotation
      stringTunings = \stringTuning <do, fa, la, re sol do'>
    }
  indent = #0
  \context {
    \Score
    supportNonIntegerFret = ##t
  }
}

<<
  \new TabStaff <<\notes>>
>>
\end{lilypond}
\end{center}


\newexercise{}


\begin{center}
\begin{lilypond}
\include "arabic.ly"
notes = \relative {
fa,4 sol' re sol fa, la fa sol' re sol re sol do sol do sol fa, la fa sol' fa,2 r2
\bar "||"
}

\layout {
  \omit Voice.StringNumber
  \context {
      \TabStaff
      \tabFullNotation
      stringTunings = \stringTuning <do, fa, la, re sol do'>
    }
  indent = #0
  \context {
    \Score
    supportNonIntegerFret = ##t
  }
}

<<
  \new TabStaff <<\notes>>
>>
\end{lilypond}
\end{center}


\newexercise{}


\begin{center}
\begin{lilypond}
\include "arabic.ly"
notes = \relative {
sol4 la sib la sol sib la sol sol la sib la sol sib la sol la2 re, sol fa,
\bar "||"
}

\layout {
  \omit Voice.StringNumber
  \context {
      \TabStaff
      \tabFullNotation
      stringTunings = \stringTuning <do, fa, la, re sol do'>
    }
  indent = #0
  \context {
    \Score
    supportNonIntegerFret = ##t
  }
}

<<
  \new TabStaff <<\notes>>
>>
\end{lilypond}
\end{center}


\newexercise{}


\begin{center}
\begin{lilypond}
\include "arabic.ly"
notes = \relative {
re4 mi fa re 
re fa mi re 
re mi fa re 
re fa mi re 
mi2 la, re r2
\bar "||"
}

\layout {
  \omit Voice.StringNumber
  \context {
      \TabStaff
      \tabFullNotation
      stringTunings = \stringTuning <do, fa, la, re sol do'>
    }
  indent = #0
  \context {
    \Score
    supportNonIntegerFret = ##t
  }
}

<<
  \new TabStaff <<\notes>>
>>
\end{lilypond}
\end{center}


\newexercise{}


\begin{center}
\begin{lilypond}
\include "arabic.ly"
notes = \relative {
re4 mi fa re 
re fa mi re 
re2 sol 
sol4 la sib la
sol sib la sol
la re, sol la,
\bar "||"
}

\layout {
  \omit Voice.StringNumber
  \context {
      \TabStaff
      \tabFullNotation
      stringTunings = \stringTuning <do, fa, la, re sol do'>
    }
  indent = #0
  \context {
    \Score
    supportNonIntegerFret = ##t
  }
}

<<
  \new TabStaff <<\notes>>
>>
\end{lilypond}
\end{center}


\newexercise{}


\begin{center}
\begin{lilypond}
\include "arabic.ly"
notes = \relative {
re4 mi fa sol la sib do re do sib la sol fa mi re mi fa la do la fa2 r2
\bar "||"
}

\layout {
  \omit Voice.StringNumber
  \context {
      \TabStaff
      \tabFullNotation
      stringTunings = \stringTuning <do, fa, la, re sol do'>
    }
  indent = #0
  \context {
    \Score
    supportNonIntegerFret = ##t
  }
}

<<
  \new TabStaff <<\notes>>
>>
\end{lilypond}
\end{center}


\newexercise{}


\begin{center}
\begin{lilypond}
\include "arabic.ly"
notes = \relative {
re4 fa re sol fa la sol sib la do do la sib sol la fa 
sol mi fa re re fa mi sol fa la sol mi fa2r2
\bar "||"
}

\layout {
  \omit Voice.StringNumber
  \context {
      \TabStaff
      \tabFullNotation
      stringTunings = \stringTuning <do, fa, la, re sol do'>
    }
  indent = #0
  \context {
    \Score
    supportNonIntegerFret = ##t
  }
}

<<
  \new TabStaff <<\notes>>
>>
\end{lilypond}
\end{center}


\newexercise{}


\begin{center}
\begin{lilypond}
\include "arabic.ly"
notes = \relative {
sol8 la\upbow sib la\upbow 
sol fa\upbow fa sol\upbow
fa mi re mi fa4\downbow la8\downbow  do\upbow
sib8 la sol la sol fa mi fa sol re mi sol fa4 r4
\bar "||"
}

\layout {
  \omit Voice.StringNumber
  \context {
      \TabStaff
      \tabFullNotation
      stringTunings = \stringTuning <do, fa, la, re sol do'>
    }
  indent = #0
  \context {
    \Score
    supportNonIntegerFret = ##t
  }
}

<<
  \new TabStaff <<\notes>>
>>
\end{lilypond}
\end{center}


\newexercise{}


\begin{center}
\begin{lilypond}
\include "arabic.ly"
notes = \relative {
re8 fa mi sol fa la sol sib la do do la sib4 la
do8 la sib sol la fa sol mi fa la do la sol4 fa
\bar "||"
}

\layout {
  \omit Voice.StringNumber
  \context {
      \TabStaff
      \tabFullNotation
      stringTunings = \stringTuning <do, fa, la, re sol do'>
    }
  indent = #0
  \context {
    \Score
    supportNonIntegerFret = ##t
  }
}

<<
  \new TabStaff <<\notes>>
>>
\end{lilypond}
\end{center}


\newexercise{}


\begin{center}
\begin{lilypond}
\include "arabic.ly"
notes = \relative {
sol4 la si la sol si la sol re mi fad mi re fad mi re sol la si la sol si la sol 
la2 re, sol sol,
\bar "||"
}

\layout {
  \omit Voice.StringNumber
  \context {
      \TabStaff
      \tabFullNotation
      stringTunings = \stringTuning <do, fa, la, re sol do'>
    }
  indent = #0
  \context {
    \Score
    supportNonIntegerFret = ##t
  }
}

<<
  \new TabStaff <<\notes>>
>>
\end{lilypond}
\end{center}


\newexercise{}


\begin{center}
\begin{lilypond}
\include "arabic.ly"
notes = \relative {
sol4 la si do si la sol fad mire mi fad sol re sol2
\bar "||"
}

\layout {
  \omit Voice.StringNumber
  \context {
      \TabStaff
      \tabFullNotation
      stringTunings = \stringTuning <do, fa, la, re sol do'>
    }
  indent = #0
  \context {
    \Score
    supportNonIntegerFret = ##t
  }
}

<<
  \new TabStaff <<\notes>>
>>
\end{lilypond}
\end{center}


\newexercise{}


\begin{center}
\begin{lilypond}
\include "arabic.ly"
notes = \relative {
sol si la do si re do la si sol la fad 
sol mi fad re re fad mi sol fad la sol 
si la re, mi fad
sol2 sol,
\bar "||"
}

\layout {
  \omit Voice.StringNumber
  \context {
      \TabStaff
      \tabFullNotation
      stringTunings = \stringTuning <do, fa, la, re sol do'>
    }
  indent = #0
  \context {
    \Score
    supportNonIntegerFret = ##t
  }
}

<<
  \new TabStaff <<\notes>>
>>
\end{lilypond}
\end{center}




\section{Ejercicios técnicos}

Los ejercicios de esta sección trabajan las principales dificultades técnicas del ud. Buena parte de ellos están relacionadas con las formas de ornamentacion melódica que vimos en el capítulo \ref{ornamentacion}.

\subsection{Ejercicios de trémolo}

\newexercise{}


\begin{center}
\begin{lilypond}
\include "arabic.ly"
notes = \relative {
\key sol \hijaz
do8 do8:32 re8 re8:32 mib8 mib8:32 fa8 fa8:32 sol8 sol8:32 lab8 lab8:32 si8 si8:32 r4
si4:32 lab4:32 sol4:32 fa4:32 lab4:32 sol4:32 fa4:32 mib4:32 sol4:32 fa4:32 mib4:32 re4:32 do1
\bar "||"
}

\layout {
  \omit Voice.StringNumber
  \context {
      \TabStaff
      \tabFullNotation
      stringTunings = \stringTuning <do, fa, la, re sol do'>
    }
  indent = #0
  \context {
    \Score
    supportNonIntegerFret = ##t
  }
}

<<
  \new TabStaff <<\notes>>
>>
\end{lilypond}
\end{center}


\audioclip{Ejercicio \arabic{exercise}.}

\newexercise{}


\begin{center}
\begin{lilypond}
\include "arabic.ly"
notes = \relative {
\key do \rast
\time 3/4
\tuplet 5/4 {do16 do' do do do}
\tuplet 5/4 {re, do' do do do}
\tuplet 5/4 {misb,16 do' do do do}
\tuplet 5/4 {re,16 do'4:32}
\tuplet 5/4 {misb,16 do'4:32}
\tuplet 5/4 {fa,16 do'4:32}
\tuplet 5/4 {misb,16 do'4:32}
\tuplet 5/4 {fa,16 do'4:32}
\tuplet 5/4 {sol16 do4:32}
\tuplet 5/4 {fa,16 do'4:32}
\tuplet 5/4 {sol16 do4:32}
\tuplet 5/4 {la16 do4:32}
\tuplet 5/4 {sol16 do4:32}
\tuplet 5/4 {la16 do4:32}
\tuplet 5/4 {sisb16 do4:32}
do2.
\tuplet 5/4 {do16 do4:32}
\tuplet 5/4 {sisb16 do4:32}
\tuplet 5/4 {la16 do4:32}
\tuplet 5/4 {sisb16 do4:32}
\tuplet 5/4 {la16 do4:32}
\tuplet 5/4 {sol16 do4:32}
\tuplet 5/4 {la16 do4:32}
\tuplet 5/4 {sol16 do4:32}
\tuplet 5/4 {fa,16 do'4:32}
\tuplet 5/4 {sol16 do4:32}
\tuplet 5/4 {fa,16 do'4:32}
\tuplet 5/4 {misb,16 do'4:32}
\tuplet 5/4 {fa,16 do'4:32}
\tuplet 5/4 {misb,16 do'4:32}
\tuplet 5/4 {re,16 do'4:32}
do,2.
\bar "||"
}

\layout {
  \omit Voice.StringNumber
  \context {
      \TabStaff
      \tabFullNotation
      stringTunings = \stringTuning <do, fa, la, re sol do'>
    }
  indent = #0
  \context {
    \Score
    supportNonIntegerFret = ##t
  }
}

<<
  \new TabStaff <<\notes>>
>>
\end{lilypond}
\end{center}


\audioclip{Ejercicio \arabic{exercise}.}

\subsection{Ejercicios de \emph{glissando}}

El aspecto más importante de estos ejercicios es prestar atención a la entonación, en especial la de la segunda nota de las dos que se conectan mediante el \emph{glissando}.

La partitura contiene la digitación a emplear en cada caso.

\newexercise{}


\begin{center}
\begin{lilypond}
\include "arabic.ly"
notes = \relative {
\override Glissando.breakable = ##t
\override Glissando.after-line-breaking = ##t
\key sol \hijaz
\time 4/4
\partial 4. do8 re mib\glissando
fa4.-1 mib16 re do r16 do8 re mib\glissando
fa4.-1 mib16 re do r16 re8 mib fa\glissando
sol4.-3\3 fa16 mib re r16 re8 mib fa\glissando
sol4.-3\3 fa16 mib re r16 fa8 sol lab\glissando
si4.-1 lab16 sol fa r16 fa8 sol lab\glissando
si4.-1 lab16 sol fa r16 sol8 lab si\glissando
do4.-3\2 sib16 lab sol r16 sol8 lab si\glissando
do4.-3\2 sib16 lab sol r16 do,8 re mib\glissando
fa4.-1 mib16 re do r16 do8 re mib\glissando
fa4.-1 mib16 re do r16 r4.\glissando
\bar "||"
}

\layout {
  \omit Voice.StringNumber
  \context {
      \TabStaff
      \tabFullNotation
      stringTunings = \stringTuning <do, fa, la, re sol do'>
    }
  indent = #0
  \context {
    \Score
    supportNonIntegerFret = ##t
  }
}

<<
  \new TabStaff <<\notes>>
>>
\end{lilypond}
\end{center}


\audioclip{Ejercicio \arabic{exercise}.}

\subsection{Ejercicios de tresillos}

Estos ejercicios están pensados para practicar la pulsación en tresillos según dos patrones: descendente--ascendente--descendente (\downbow \upbow \downbow) y descendente--descendente--ascendente (\downbow \downbow \upbow). El patron correspondiente se indica en cada caso. Es importante prestar atención a los cambios de cuerda, donde reside en muchos casos la mayor dificultad.

\newexercise{}


\begin{center}
\begin{lilypond}
\include "arabic.ly"
notes = \relative {
\override Stem.direction = #down
\tuplet 3/2 {la,8\downbow re\downbow re\upbow}
\tuplet 3/2 {la\downbow re\downbow re\upbow}
\tuplet 3/2 {re sol sol}
\tuplet 3/2 {re sol sol}
\tuplet 3/2 {sol do do}
\tuplet 3/2 {sol do do}
\tuplet 3/2 {sol do do}
\tuplet 3/2 {sol do do}
\tuplet 3/2 {re, sol sol}
\tuplet 3/2 {re sol sol}
\tuplet 3/2 {la, re re}
\tuplet 3/2 {la re re}
\bar "||"
}

\layout {
  \omit Voice.StringNumber
  \context {
      \TabStaff
      \tabFullNotation
      stringTunings = \stringTuning <do, fa, la, re sol do'>
    }
  indent = #0
  \context {
    \Score
    supportNonIntegerFret = ##t
  }
}

<<
  \new TabStaff <<\notes>>
>>
\end{lilypond}
\end{center}


\audioclip{Ejercicio \arabic{exercise}.}

\newexercise{}


\begin{center}
\begin{lilypond}
\include "arabic.ly"
notes = \relative {
\override Stem.direction = #down
\tuplet 3/2 {do'8\downbow sol\downbow sol\upbow}
\tuplet 3/2 {do\downbow sol\downbow sol\upbow}
\tuplet 3/2 {sol re re}
\tuplet 3/2 {sol re re}
\tuplet 3/2 {re la la}
\tuplet 3/2 {re la la}
\tuplet 3/2 {re la la}
\tuplet 3/2 {re la la}
\tuplet 3/2 {sol' re re}
\tuplet 3/2 {sol re re}
\tuplet 3/2 {do' sol sol}
\tuplet 3/2 {do sol sol}
\bar "||"
}

\layout {
  \omit Voice.StringNumber
  \context {
      \TabStaff
      \tabFullNotation
      stringTunings = \stringTuning <do, fa, la, re sol do'>
    }
  indent = #0
  \context {
    \Score
    supportNonIntegerFret = ##t
  }
}

<<
  \new TabStaff <<\notes>>
>>
\end{lilypond}
\end{center}


\audioclip{Ejercicio \arabic{exercise}.}

\newexercise{}

Recorrer un \maqam alternando sus notas con la doble pulsación de una de ellas (en este caso la nota sol, \emph{ghammaz} de la escala que se recorre). Este ejercicio puede adaptarse a otros \maqamaatb. También puede variarse la nota repetida, utilizando por ejemplo el do de la octava superior, de manera que se aumenta la distancia y el salto de cuerdas necesario.


\begin{center}
\begin{lilypond}
\include "arabic.ly"
notes = \relative {
\key do \rast
\override Stem.direction = #down
\tuplet 3/2 {do8\downbow sol'\downbow sol\upbow}
\tuplet 3/2 {re\downbow sol\downbow sol\upbow}
\tuplet 3/2 {mib sol sol}
\tuplet 3/2 {fa sol sol}
\tuplet 3/2 {sol sol sol}
\tuplet 3/2 {lab sol sol}
\tuplet 3/2 {si sol sol}
\tuplet 3/2 {do sol sol}
\bar "||"
}

\layout {
  \omit Voice.StringNumber
  \context {
      \TabStaff
      \tabFullNotation
      stringTunings = \stringTuning <do, fa, la, re sol do'>
    }
  indent = #0
  \context {
    \Score
    supportNonIntegerFret = ##t
  }
}

<<
  \new TabStaff <<\notes>>
>>
\end{lilypond}
\end{center}


\audioclip{Ejercicio \arabic{exercise}.}

\newexercise{}


\begin{center}
\begin{lilypond}
\include "arabic.ly"
notes = \relative {
\key do \minor
\time 3/4
\override Stem.direction = #down
\tuplet 3/2 {do8\downbow do\upbow do\downbow}
\tuplet 3/2 {re\downbow re\upbow re\downbow}
\tuplet 3/2 {mib mib mib}
\tuplet 3/2 {fa fa fa}
\tuplet 3/2 {sol sol sol}
\tuplet 3/2 {la la la}
\tuplet 3/2 {sib sib sib}
\tuplet 3/2 {do do do}
\tuplet 3/2 {re re re}
mib4 r2
\tuplet 3/2 {mib8 mib mib}
\tuplet 3/2 {re re re}
\tuplet 3/2 {do do do}
\tuplet 3/2 {sib sib sib}
\tuplet 3/2 {la la la}
\tuplet 3/2 {sol sol sol}
\tuplet 3/2 {fa fa fa}
\tuplet 3/2 {mib mib mib}
\tuplet 3/2 {re re re}
do4 r2
\bar "||"
}

\layout {
  \omit Voice.StringNumber
  \context {
      \TabStaff
      \tabFullNotation
      stringTunings = \stringTuning <do, fa, la, re sol do'>
    }
  indent = #0
  \context {
    \Score
    supportNonIntegerFret = ##t
  }
}

<<
  \new TabStaff <<\notes>>
>>
\end{lilypond}
\end{center}


\audioclip{Ejercicio \arabic{exercise}.}

\section{Ejercicios de independencia de dedos}

Algunos ejercicios cromáticos para ejercitar la independencia de los dedos de la mano izquierda.

%https://appliedguitartheory.com/lessons/10-finger-independence-exercises-guitar/

\newexercise{}


\section{Ejercicios rítmicos}

Los siguientes ejercicios se basan en los principales ritmos (\emph{iqaa}) árabes, presentando su estructura principal y algunas variaciones rítmicas.

\newexercise{Ritmo \emph{malfuf}}


\begin{center}
\begin{lilypond}
\include "arabic.ly"
notes = \relative {
\key do \minor
\time 3,3,2 2/4
do8. mib8. mib8
do8. mib8. mib8
do8. mib8. mib8
do8. mib8. mib8
do8. mib8. mib16 (re)
do8. mib8. mib8
do8. mib8. mib16 (re)
do8. mib8. mib8
do16-> do16 do16 mib16-> mib16 mib16 mib8->
do8. mib8. mib8
do16-> do16 do16 mib16-> mib16 mib16 mib8->
do8. mib8. mib8
do16-> do16 do16 mib16-> mib16 mib16 mib16-> (re)
do8. mib8. mib8
do16-> do16 do16 mib16-> mib16 mib16 mib16-> (re)
do8. mib8. mib8
do16 sol' sol sol lab (sol) fa mib
do8. mib8. mib8
do16 sol' sol sol lab (sol) fa mib
do8. mib8. mib8
do16 sol' sol sol lab (sol) fa mib
do16 sol' sol sol lab (sol) fa mib
do16 sol' sol sol lab (sol) fa mib
do8. mib8. mib8
do8. mib8. mib8
\bar "||"
}

\layout {
  \omit Voice.StringNumber
  \context {
      \TabStaff
      \tabFullNotation
      stringTunings = \stringTuning <do, fa, la, re sol do'>
    }
  indent = #0
  \context {
    \Score
    supportNonIntegerFret = ##t
  }
}

<<
  \new TabStaff <<\notes>>
>>
\end{lilypond}
\end{center}


\audioclip{Ejercicio \arabic{exercise}.}

\newexercise{Ritmo \emph{maqsum}}


\begin{center}
\begin{lilypond}
\include "arabic.ly"
notes = \relative {
\time 4/4
\key do \minor
re8 sol r8 sol fa r8 sol r8
re8 sol r8 sol fa fa sol r8
re8 sol r8 sol fa r8 sol r8
re8 sol r8 sol fa fa sol r8
re8 sol sol16 sol sol sol fa8 r8 sol r8
re8 sol sol16 sol sol sol fa8 fa sol r8
re8 sol sol16 sol sol sol fa8 r8 sol r8
re8 sol sol16 sol sol sol la16 (sol) fa8 sol r8
re8 sol r8 sol fa r8 sol r8
la'16 (sib) sol (la) fa (sol) mib (fa) re (mib) do (re) sib (do) la (sib)
sol8 sol' r8 sol fa r8 sol r8
re8 sol r8 sol fa fa sol r8
re8 sol r8 sol fa r8 sol r8
la'16 (sib) sol (la) fa (sol) mib (fa) re (mib) do (re) sib (do) la (sib)
sol4 r2.
\bar "||"
}

\layout {
  \omit Voice.StringNumber
  \context {
      \TabStaff
      \tabFullNotation
      stringTunings = \stringTuning <do, fa, la, re sol do'>
    }
  indent = #0
  \context {
    \Score
    supportNonIntegerFret = ##t
  }
}

<<
  \new TabStaff <<\notes>>
>>
\end{lilypond}
\end{center}


\audioclip{Ejercicio \arabic{exercise}.}

\newexercise{Ritmo \emph{baladi}}


\begin{center}
\begin{lilypond}
\include "arabic.ly"
notes = \relative {
\time 4/4
\key do \rast
re8 re r8 do re r8 do r8
re8 re r8 do re r8 do r8
re8 re do16 do do8 re do16 do do8 do16 do
re8 re do16 do do8 re do16 do do8 do16 do
re8 re do16 do do8 re misb16 fa misb re do8
re8 re do16 do do8 re misb16 fa misb re do8
re8 re do16 do do8 re fa8:64 misb16 re do8
re8 re do16 do do8 re fa8:64 misb16 re do8
\bar "||"
}

\layout {
  \omit Voice.StringNumber
  \context {
      \TabStaff
      \tabFullNotation
      stringTunings = \stringTuning <do, fa, la, re sol do'>
    }
  indent = #0
  \context {
    \Score
    supportNonIntegerFret = ##t
  }
}

<<
  \new TabStaff <<\notes>>
>>
\end{lilypond}
\end{center}


\audioclip{Ejercicio \arabic{exercise}.}

\section{Ejercicios por \maqamb}

\subsection{\emph{Maqam} Rast}
\subsection{\emph{Maqam} Bayati}
\subsection{\emph{Maqam} Ayam}
\subsection{\emph{Maqam} Nahawand}

%https://issuu.com/theschoolofarmenianoud/docs/i_girq__lriv